\documentclass{article}
\usepackage{amsmath}
\usepackage{amssymb}
\usepackage{amsthm}
\usepackage{physics}
\usepackage{graphicx}

\graphicspath{ {img/} }


\newtheorem*{conjecture}{Conjecture}
\newtheorem*{theorem}{Theorem}
\newtheorem{lemma}{Lemma}
\newtheorem*{lemma*}{Lemma}
\newtheorem*{proposition}{Proposition}


\makeatletter
\newcommand*{\rom}[1]{\expandafter\@slowromancap\romannumeral #1@}
\makeatother


\begin{document}

\section*{2.1}

Let us indicate the displacement of the \(P\) particle as:
\begin{equation*}
    x(t) = 6t^{2} - t^{3} + 1
\end{equation*}
where \(x\) is measured in metres and \(t\) in seconds. The velocity and acceleration can be found
by calculcating the first and second derivative respectively:

\begin{equation*}
    v(t) = \dv{x}{t} = 12t - 3t^{2} = -3t^{2} + 12t
\end{equation*}
\begin{equation*}
    a(t) = \dv{v}{t} = -6t + 12
\end{equation*}


\(P\) is at rest when its velocity equals zero. We immediately set up a necessary equation:
\begin{equation*}
    v(t) = 0 \implies -3t^{2} + 12t = 0
\end{equation*}
\begin{equation*}
    3t(4 - t) = 0
\end{equation*}

We have found two times at which \(P\) is at rest: \(t_1 = 0s\) and \(t_2 = 4s\). We will conclude with
calculcating the position of \(P\) at these times:
\begin{equation*}
    x(0) = 6 \cdot 0^2 - 0^{3} + 1 = 1m
\end{equation*}
\begin{equation*}
    x(4) = 6 \cdot 4^{2} - 4^{3} + 1 = 96 - 64 + 1 = 33m
\end{equation*} 

\end{document}